% ----------------------------------------------------------------------
%  3.4.4. Анализ Даркена
% ----------------------------------------------------------------------
\subsection{Анализ Даркена в термодинамической теории диффузии}\label{ssec:darken}

\paragraph{Эффект Киркендалла}
Смигельскас и Киркендалл поставили простой опыт
(рис.~\ref{fig:kirkendall}). На поверхность латуни (около 30~\% Zn)
положили тонкие молибденовые проволочки-метки и нарастили сверху слой
меди. После длительного отжига метки, не участвующие в диффузии,
\emph{сместились в сторону латуни}. Вывод: компоненты сплава
диффундируют с разной скоростью. Цинк уходит из латуни в медь быстрее,
чем медь приходит в латунь; латунь «худеет», и метки движутся вместе с
решёткой в её сторону. Это и есть \term{эффект Киркендалла}.

\begin{figure}[ht]
\centering
\begin{tikzpicture}[x=1cm, y=1cm]
  \foreach \sh/\m/\lab in {0/0/{\textit{а}~— до отжига}, 7.2/0.55/{\textit{б}~— после отжига}}{
    \begin{scope}[shift={(\sh,0)}]
      \fill[fmAmber!25] (0,0) rectangle (3,2.4);
      \fill[fmAmber!55] (3,0) rectangle (6,2.4);
      \ifdim\m pt>0pt
        \shade[left color=fmAmber!25, right color=fmAmber!55] (2.2,0) rectangle (4.2,2.4);
        \foreach \p in {(3.95,0.45),(4.25,1.4),(3.85,1.95),(4.4,0.8),(4.1,1.05),(4.5,1.75)}{\fill[white] \p circle (0.07);}
      \fi
      \draw[atomLine] (0,0) rectangle (6,2.4);
      \foreach \y in {0.4,1.2,2.0}{\fill[fmBlue] ({3+\m},\y) circle (0.09);}
      \draw[fmBlue, dashed] ({3+\m},0) -- ({3+\m},2.4);
      \node[font=\small] at (1.5,2.75) {Cu};
      \node[font=\small] at (4.5,2.75) {Cu + 30\,\% Zn};
      \node[font=\small] at (3,-0.45) {\lab};
    \end{scope}}
  \draw[-{Stealth[length=2mm]}, fmRed, thick] (11.1,3.25) -- ++(-1.0,0);
  \draw[-{Stealth[length=2mm]}, fmBlue, thick] (9.75,-0.95) -- ++(0.5,0);
  \node[font=\scriptsize, fmRed, anchor=west] at (11.15,3.25) {Zn};
  \node[font=\scriptsize, fmBlue, anchor=east] at (9.7,-0.95) {Cu};
\end{tikzpicture}
\caption{Опыт Смигельскаса и Киркендалла (схема): \textit{а}~— метки
(синие точки) на исходной границе медь~— латунь; \textit{б}~— после
отжига метки сместились в сторону латуни, а в латуни у зоны диффузии
появились поры (эффект Френкеля). Поток цинка больше встречного потока
меди}
\label{fig:kirkendall}
\end{figure}

\paragraph{Уравнения Даркена}
Даркен описал эффект количественно. Каждому компоненту он приписал
собственный (\emph{парциальный}) коэффициент диффузии $D_1$, $D_2$,
задающий поток относительно решётки (меток). Сама решётка движется
относительно неподвижной системы координат со скоростью $v$.
\begin{shagi}
\item Потоки в неподвижной системе:
$J_i'=-D_i\,\partial C_i/\partial x+vC_i$. При постоянном объёме
(и пренебрегая вакансиями) $J_1'=-J_2'$.
\item Отсюда скорость меток
\end{shagi}
\begin{keyeq}\label{eq:darken-v}
v=(D_1-D_2)\frac{\partial N_1}{\partial x},
\end{keyeq}
где $N_i=C_i/C$~— мольные доли.
\begin{shagi}[resume]
\item Коэффициент \term{взаимной диффузии} $\widetilde D$ определим как
коэффициент пропорциональности между полным потоком любого компонента
и его градиентом. Подставив $v$, получим
\end{shagi}
\begin{keyeq}\label{eq:darken-D}
\widetilde D=N_2D_1+N_1D_2 .
\end{keyeq}

$\widetilde D$~— мера скорости \emph{гомогенизации} сплава: он описывает
поток через плоскость, сквозь которую в противоположных направлениях
проходит равное число атомов (плоскость Матано). Измерив $\widetilde D$
и скорость меток $v$, по \eqref{eq:darken-v} и \eqref{eq:darken-D}
находят оба парциальных коэффициента. Поскольку $D_1\neq D_2$,
$\widetilde D$ неизбежно зависит от состава. Правда, опыт показывает, что
допущения Даркена выполняются не всегда даже приближённо: мешают
пористость и изменение поперечных размеров образца в зоне диффузии.

\begin{primer}[Скорость меток]
Пусть в некоторой плоскости $N_1=0{,}3$, $D_1=3\cdot10^{-14}$~м$^2$/с,
$D_2=10^{-14}$~м$^2$/с, а $\partial N_1/\partial x=-200$~м$^{-1}$ (состав
меняется на 0,2 на миллиметре). Тогда
$\widetilde D=0{,}7\cdot3\cdot10^{-14}+0{,}3\cdot10^{-14}=2{,}4\cdot10^{-14}$~м$^2$/с,
а $v=(3-1)\cdot10^{-14}\cdot(-200)=-4\cdot10^{-12}$~м/с: метки движутся в
сторону, где больше быстрого компонента, со скоростью около
0,35~мкм в сутки.
\end{primer}

\paragraph{Термодинамический множитель}
Поток пропорционален градиенту химического потенциала (п.~\ref{ssec:onsager}):
$J_i=-L_{ii}\operatorname{grad}\mu_i$. С $\mu_i=\mu_i^0+RT\ln(f_iC_i)$
($f_i$~— коэффициент активности) получаем
\begin{equation}\label{eq:thermo-factor}
D_i=\frac{L_{ii}}{C_i}\frac{\partial\mu_i}{\partial\ln C_i}
=RT\frac{L_{ii}}{C_i}\left(1+\frac{\partial\ln f_i}{\partial\ln C_i}\right).
\end{equation}
Даркен предположил, что \emph{подвижность} $L_{ii}/C_i$ при само- и
взаимной диффузии одинакова. Тогда $RTL_{ii}/C_i=D_i^*$~— коэффициент
самодиффузии (меченых атомов), и
\begin{keyeq}\label{eq:darken-full}
D_i=D_i^*\left(1+\frac{\partial\ln f_i}{\partial\ln N_i}\right),\qquad
\widetilde D=\bigl(N_1D_2^*+N_2D_1^*\bigr)\left(1+\frac{\partial\ln f_1}{\partial\ln N_1}\right).
\end{keyeq}
Выражение в скобках~— \term{термодинамический множитель}; для идеального
раствора он равен единице. (Для двух компонентов он одинаков: по
уравнению Гиббса--Дюгема $\partial\ln f_1/\partial\ln N_1=\partial\ln f_2/\partial\ln N_2$.)
Формула \eqref{eq:darken-full} связывает самодиффузию~— теоретически
самую простую величину (среднее число беспорядочных скачков атома в
секунду в однородном сплаве)~— со скоростью взаимной диффузии. Опыты
многих авторов её подтвердили.

\paragraph{Роль вакансий: поправка Маннинга}
Даркен не учёл вакансий. Но раз потоки атомов не равны, через решётку
идёт \emph{результирующий поток вакансий}, $J_v\propto(D_1^*-D_2^*)\,\partial C_1/\partial x$.
Ле Клер показал, что этот поток целиком определяется недиагональными
коэффициентами Онзагера, связывающими вакансии с атомами. Маннинг
заметил, что поток вакансий делает скачки направленно неравноправными:
против потока вакансий атомы прыгают чаще. Поправка увеличивает
$D_1$ быстрого компонента и уменьшает $D_2$:
\[
D_i=D_i^*\Bigl(1+\frac{\partial\ln f_i}{\partial\ln C_i}\Bigr)(1\pm v_i),\qquad
\widetilde D=\bigl(C_1D_2^*+C_2D_1^*\bigr)\Bigl(1+\frac{\partial\ln f_1}{\partial\ln C_1}\Bigr)S,
\]
\[
S=1+\frac{2C_1C_2(D_1^*-D_2^*)^2}{M_0(C_1D_1^*+C_2D_2^*)(C_1D_2^*+C_2D_1^*)},
\]
где $M_0$ зависит от решётки: 7,15 (ГЦК), 5,33 (ОЦК), 3,75 (простая
кубическая), 2 (алмаз). Множитель $S$ максимален при равных
концентрациях и растёт с $\beta=D_1^*/D_2^*$ (рис.~\ref{fig:manning}), но
даже при $\beta\to\infty$ для ГЦК не превышает $1+2/M_0=1{,}28$; обычно
$\beta\le3$ и $S\approx1{,}05$~— в пределах точности опыта.

\begin{figure}[ht]
\centering
\begin{tikzpicture}
\begin{axis}[width=0.58\textwidth, height=5cm, xmin=0, xmax=1, ymin=1, ymax=1.32,
  xlabel={$C_1$}, ylabel={$S$}, axis lines=left,
  yticklabel style={/pgf/number format/fixed, /pgf/number format/precision=2, /pgf/number format/use comma}]
  \addplot[fmBlue, thick, domain=0:1, samples=60] {1+2*x*(1-x)*(3-1)^2/(7.15*(3*x+(1-x))*(x+3*(1-x)))};
  \addplot[fmGreen!60!black, thick, domain=0:1, samples=60] {1+2*x*(1-x)*(10-1)^2/(7.15*(10*x+(1-x))*(x+10*(1-x)))};
  \addplot[fmRed, thick, domain=0:1, samples=60] {1+2*x*(1-x)*(1000-1)^2/(7.15*(1000*x+(1-x))*(x+1000*(1-x)))};
  \node[font=\scriptsize, fmBlue] at (axis cs:0.5,1.065) {$\beta=3$};
  \node[font=\scriptsize, fmGreen!50!black] at (axis cs:0.5,1.16) {$\beta=10$};
  \node[font=\scriptsize, fmRed] at (axis cs:0.5,1.295) {$\beta\to\infty$};
\end{axis}
\end{tikzpicture}
\caption{Поправочный множитель Маннинга к формуле Даркена для ГЦК-сплавов
(расчёт): максимальная поправка~— 28~\% при $\beta\to\infty$}
\label{fig:manning}
\end{figure}

\paragraph{Эффекты Киркендалла и Френкеля}
Поток вакансий в зоне диффузии неоднороден: $\partial C_1/\partial x$
меняется от нуля до максимума. Чтобы концентрация вакансий оставалась
равновесной, в одной части зоны они должны рождаться (источники~—
дислокации), в другой~— исчезать на стоках или собираться в поры. Изменение
объёма разных частей зоны и даёт сдвиг меток (эффект Киркендалла), а
скопление пор~— \term{эффект Френкеля} (рис.~\ref{fig:kirkendall},\,\textit{б}).
Пористость ослабляет соединения разнородных металлов, поэтому эти
эффекты важны для паяных и сварных узлов и для взаимодействия топлива
с оболочкой твэла.

\begin{podrobnee}[многокомпонентные системы]
Для $n$ компонентов (без перекрёстных коэффициентов, но с эффектом
Киркендалла) состав задают $n-1$ концентраций, и
\[
\frac{\partial C_i}{\partial t}=\frac{\partial}{\partial x}\sum_{k=1}^{n-1}\widetilde D_{ik}\frac{\partial C_k}{\partial x},
\qquad i=1,\ldots,n-1 ,
\]
где $\widetilde D_{ik}$~— матрица коэффициентов взаимной диффузии. Её
элементы выражаются через коэффициенты самодиффузии $D_m^*$ и
термодинамические множители $m_{ik}=\delta_{ik}+\partial\ln f_i/\partial\ln C_k$
(обобщённое уравнение Даркена):
$\widetilde D_{ik}=\sum_{m=1}^{n-1}\bigl[\delta_{im}D_m^*-C_i(D_m^*-D_n^*)\bigr]\frac{C_m}{C_k}m_{mk}$.
Для тройной системы непосредственно определяются четыре коэффициента.
Точных методов решения этой системы нет.
\end{podrobnee}
